% ============================================================
\chapter{Classical Text Representations}
\label{ch:representations}
% ============================================================

Here is the fundamental problem of natural language processing: computers operate on numbers, but humans communicate with words. Bridging this gap---turning text into mathematical objects that algorithms can manipulate---is the challenge of \emph{text representation}, and the history of NLP can be read as a series of increasingly sophisticated answers to this single question.

The earliest approach, dating back to the 1950s, was brutally simple: represent each document as a vector of word counts, ignoring grammar, word order, and meaning entirely. This ``bag of words'' model is laughably crude, yet it powered information retrieval systems for decades. The refinement came with TF-IDF weighting, which recognised that not all words are equally informative: ``the'' appears everywhere and tells you nothing, while ``mitochondria'' is rare and highly diagnostic. From there, the quest for richer representations led to latent semantic analysis, topic models, and ultimately to the word embeddings of Mikolov's Word2Vec (2013)---the breakthrough that showed meaning could be captured by geometry. But before we reach those heights, we must understand the classical foundations.

\begin{intuition}[title=\textbf{From text to vectors}]
Machine learning algorithms operate on numerical vectors, not raw text. This
chapter presents classical methods for transforming textual documents into
vector representations that can be used by classifiers or search engines.
\end{intuition}

% ------------------------------------------------------------
\section{Bag-of-Words Representation}
\label{sec:rep:bow}
% ------------------------------------------------------------

\begin{definition}[Bag of Words]
The \emph{bag-of-words} (BoW) model represents a document $d$ as a vector
$\mathbf{x}_d \in \N^V$ where the $j$-th component is the count of term $t_j$
in $d$:
\[
x_{d,j} = \text{count}(t_j, d)
\]
\end{definition}

This model completely ignores word order. The sentences ``the cat eats the
mouse'' and ``the mouse eats the cat'' have the same representation.

\begin{remark}
Despite its simplicity, the bag-of-words model remains useful as a baseline and
in applications where word order is less important (topic classification,
document filtering).
\end{remark}

% ------------------------------------------------------------
\section{TF-IDF Weighting}
\label{sec:rep:tfidf}
% ------------------------------------------------------------

Raw frequency does not reflect the discriminative importance of a term.
Very frequent words (``the'', ``of'', ``and'') dominate the representation
without providing semantic information.

\begin{definition}[TF -- Term Frequency]
The \emph{term frequency} of $t$ in document $d$ is:
\[
\text{tf}(t, d) = \frac{\text{count}(t, d)}{\sum_{t' \in d} \text{count}(t', d)}
\]
\end{definition}

\begin{definition}[IDF -- Inverse Document Frequency]
The \emph{inverse document frequency} of term $t$ in a corpus $\mathcal{D}$
of $N$ documents is:
\[
\text{idf}(t, \mathcal{D}) = \log \frac{N}{\abs{\{d \in \mathcal{D} : t \in d\}}}
\]
\end{definition}

\begin{definition}[TF-IDF]
The \emph{TF-IDF} weighting combines both measures:
\[
\text{tf-idf}(t, d, \mathcal{D}) = \text{tf}(t, d) \times \text{idf}(t, \mathcal{D})
\]
\end{definition}

\begin{theorem}[Properties of TF-IDF]
Let $\mathbf{v}_d \in \R^V$ be the TF-IDF vector of document $d$. Then:
\begin{enumerate}
    \item $\text{tf-idf}(t, d) = 0$ if $t \notin d$;
    \item $\text{tf-idf}(t, d)$ is high if $t$ is frequent in $d$ but rare in
          the corpus;
    \item $\text{idf}(t) = 0$ if $t$ appears in all documents.
\end{enumerate}
\end{theorem}

\begin{keyformulas}[title=\textbf{TF-IDF Variants}]
\begin{align}
\text{tf}_{\log}(t,d) &= 1 + \log(\text{count}(t,d)) & \text{(log tf)} \\
\text{idf}_{\text{smooth}}(t) &= \log\frac{N+1}{\text{df}(t)+1} + 1 & \text{(smoothed idf)} \\
\text{tf-idf}_{\text{norm}} &= \frac{\text{tf-idf}(t,d)}{\norm{\mathbf{v}_d}} & \text{($L_2$ normalized)}
\end{align}
\end{keyformulas}

\begin{codeblock}[title=\textbf{TF-IDF with scikit-learn}]
\begin{minted}{python}
from sklearn.feature_extraction.text import TfidfVectorizer

corpus = [
    "The cat sleeps on the mat",
    "The dog runs in the garden",
    "The cat and the dog play together"
]

vectorizer = TfidfVectorizer()
X = vectorizer.fit_transform(corpus)
print(f"Matrix shape: {X.shape}")
print(f"Vocabulary: {vectorizer.get_feature_names_out()}")
print(f"TF-IDF matrix (dense):\n{X.toarray().round(2)}")
\end{minted}
\end{codeblock}

% ------------------------------------------------------------
\section{$N$-gram Models}
\label{sec:rep:ngrams}
% ------------------------------------------------------------

\begin{definition}[$N$-gram]
An \emph{$n$-gram} is a contiguous subsequence of $n$ tokens extracted from a
text. For a sequence $(x_1, \ldots, x_T)$, the $n$-grams are:
\[
\{(x_t, x_{t+1}, \ldots, x_{t+n-1}) : t = 1, \ldots, T - n + 1\}
\]
\end{definition}

\subsection{$N$-gram Language Model}

An $n$-gram language model approximates the conditional probability using a
Markov assumption of order $n-1$:

\begin{definition}[$N$-gram Language Model]
\[
P(x_t \mid x_1, \ldots, x_{t-1}) \approx P(x_t \mid x_{t-n+1}, \ldots, x_{t-1})
\]
estimated by maximum likelihood:
\[
\hat{P}(x_t \mid x_{t-n+1}^{t-1}) = \frac{\text{count}(x_{t-n+1}, \ldots, x_t)}{\text{count}(x_{t-n+1}, \ldots, x_{t-1})}
\]
\end{definition}

\begin{warning}[title=\textbf{Zero probability problem}]
If an $n$-gram never appears in the training corpus, its estimated probability
is zero, which makes the perplexity infinite. Smoothing techniques (Laplace,
Kneser-Ney) are essential.
\end{warning}

\subsection{Laplace Smoothing (Add-$\alpha$)}

\begin{definition}[Laplace Smoothing]
Laplace smoothing adds a pseudo-count $\alpha > 0$:
\[
\hat{P}_\alpha(x_t \mid x_{t-n+1}^{t-1}) = \frac{\text{count}(x_{t-n+1}^t) + \alpha}{\text{count}(x_{t-n+1}^{t-1}) + \alpha V}
\]
where $V = \abs{\mathcal{V}}$.
\end{definition}

\subsection{Kneser-Ney Smoothing}

\begin{definition}[Kneser-Ney Smoothing]
Kneser-Ney smoothing uses an absolute discount $\delta$ and a continuation
distribution:
\[
P_{\text{KN}}(w \mid h) = \frac{\max(\text{count}(h, w) - \delta, 0)}{\text{count}(h)} + \lambda(h) \cdot P_{\text{cont}}(w)
\]
where $P_{\text{cont}}(w) \propto \abs{\{h' : \text{count}(h', w) > 0\}}$ is
the continuation probability and $\lambda(h)$ is a normalization factor.
\end{definition}

% ------------------------------------------------------------
\section{Document Similarity}
\label{sec:rep:similarity}
% ------------------------------------------------------------

\begin{definition}[Cosine Similarity]
The \emph{cosine similarity} between two vectors $\mathbf{u}, \mathbf{v} \in \R^V$ is:
\[
\cos(\mathbf{u}, \mathbf{v}) = \frac{\inner{\mathbf{u}}{\mathbf{v}}}{\norm{\mathbf{u}} \cdot \norm{\mathbf{v}}}
\]
\end{definition}

\begin{proposition}[Properties of Cosine Similarity]
\begin{enumerate}
    \item $\cos(\mathbf{u}, \mathbf{v}) \in [-1, 1]$
    \item $\cos(\mathbf{u}, \mathbf{v}) = 1 \iff \mathbf{u} = \lambda \mathbf{v}$
          for $\lambda > 0$
    \item For TF-IDF vectors (non-negative components):
          $\cos(\mathbf{u}, \mathbf{v}) \in [0, 1]$
\end{enumerate}
\end{proposition}

\begin{definition}[Jaccard Distance]
For two term sets $A$ and $B$:
\[
J(A, B) = \frac{\abs{A \cap B}}{\abs{A \cup B}}, \qquad d_J(A, B) = 1 - J(A, B)
\]
\end{definition}

% ------------------------------------------------------------
\section{Document-Term Matrix and SVD}
\label{sec:rep:svd}
% ------------------------------------------------------------

The document-term matrix $\mathbf{M} \in \R^{N \times V}$ is generally very
sparse ($> 99\%$ zeros). The singular value decomposition (SVD) enables
dimensionality reduction.

\begin{definition}[Latent Semantic Analysis (LSA)]
Latent Semantic Analysis (LSA) applies truncated SVD to the TF-IDF matrix:
\[
\mathbf{M} \approx \mathbf{U}_k \boldsymbol{\Sigma}_k \mathbf{V}_k^\top
\]
where $k \ll \min(N, V)$. Documents are then represented in $\R^k$ by the rows
of $\mathbf{U}_k \boldsymbol{\Sigma}_k$.
\end{definition}

\begin{theorem}[Optimality of Truncated SVD (Eckart-Young)]
The rank-$k$ matrix that minimizes the Frobenius error
$\norm{\mathbf{M} - \hat{\mathbf{M}}}_F$ is given by the truncated SVD
$\hat{\mathbf{M}} = \mathbf{U}_k \boldsymbol{\Sigma}_k \mathbf{V}_k^\top$.
\end{theorem}

\begin{codeblock}[title=\textbf{LSA with scikit-learn}]
\begin{minted}{python}
from sklearn.feature_extraction.text import TfidfVectorizer
from sklearn.decomposition import TruncatedSVD

vectorizer = TfidfVectorizer(max_features=10000)
X_tfidf = vectorizer.fit_transform(corpus)

svd = TruncatedSVD(n_components=100, random_state=42)
X_lsa = svd.fit_transform(X_tfidf)
print(f"Explained variance: {svd.explained_variance_ratio_.sum():.2%}")
print(f"Reduced dimension: {X_lsa.shape}")
\end{minted}
\end{codeblock}

% ------------------------------------------------------------
\section{Character $N$-gram Representations}
\label{sec:rep:char}
% ------------------------------------------------------------

Character $n$-grams capture morphology and are robust to spelling errors:

\begin{example}
The word ``hello'' with $n=3$ produces the trigrams:
\texttt{\#he}, \texttt{hel}, \texttt{ell}, \texttt{llo}, \texttt{lo\#}
(where \texttt{\#} marks word boundaries).
\end{example}

% ------------------------------------------------------------
\section{BM25 Model}
\label{sec:rep:bm25}
% ------------------------------------------------------------

\begin{definition}[BM25]
The BM25 score of document $d$ for query $q = \{q_1, \ldots, q_m\}$ is:
\[
\text{BM25}(d, q) = \sum_{i=1}^{m} \text{idf}(q_i) \cdot
\frac{\text{tf}(q_i, d) \cdot (k_1 + 1)}{\text{tf}(q_i, d) + k_1 \cdot \left(1 - b + b \cdot \frac{\abs{d}}{\text{avgdl}}\right)}
\]
where $k_1$ and $b$ are hyperparameters (typically $k_1 = 1.2$, $b = 0.75$)
and $\text{avgdl}$ is the average document length.
\end{definition}

\begin{bestpractice}[title=\textbf{BM25 in practice}]
BM25 remains a reference method for information retrieval. Even modern RAG
systems often use BM25 as a first retrieval stage before neural reranking.
\end{bestpractice}

% ------------------------------------------------------------
\section{Limitations of Sparse Representations}
\label{sec:rep:limits}
% ------------------------------------------------------------

The methods presented in this chapter suffer from several limitations:

\begin{enumerate}
    \item \textbf{High dimensionality}: $V$ can reach $10^5$ to $10^6$.
    \item \textbf{Sparsity}: vectors are very sparse, making distance measures
          unreliable.
    \item \textbf{Lack of semantics}: synonymous words have orthogonal
          components (``car'' $\perp$ ``automobile'').
    \item \textbf{Context independence}: each word has a fixed representation
          independent of context.
\end{enumerate}

These limitations motivate dense embeddings (Chapter~\ref{ch:embeddings}),
which represent words in a low-dimensional space where geometric proximity
reflects semantic similarity.

% ------------------------------------------------------------
\section{Application: Naive Bayes Classification}
\label{sec:rep:naivebayes}
% ------------------------------------------------------------

\begin{definition}[Multinomial Naive Bayes Classifier]
Let a document be represented by its word counts $\mathbf{x} = (x_1, \ldots, x_V)$.
The multinomial naive Bayes classifier predicts:
\[
\hat{y} = \argmax_{c \in \mathcal{Y}} \left[\log P(c) + \sum_{j=1}^{V} x_j \log P(t_j \mid c)\right]
\]
with $P(t_j \mid c) = \frac{\text{count}(t_j, c) + \alpha}{\sum_{j'} \text{count}(t_{j'}, c) + \alpha V}$.
\end{definition}

\begin{codeblock}[title=\textbf{Naive Bayes classification}]
\begin{minted}{python}
from sklearn.feature_extraction.text import TfidfVectorizer
from sklearn.naive_bayes import MultinomialNB
from sklearn.pipeline import Pipeline
from sklearn.datasets import fetch_20newsgroups

data = fetch_20newsgroups(
    subset='train',
    categories=['sci.space', 'rec.sport.baseball']
)
pipeline = Pipeline([
    ('tfidf', TfidfVectorizer(max_features=10000)),
    ('clf', MultinomialNB(alpha=0.1))
])
pipeline.fit(data.data, data.target)
print(f"Accuracy (train): {pipeline.score(data.data, data.target):.2%}")
\end{minted}
\end{codeblock}

% ------------------------------------------------------------
\section{Exercises}
\label{sec:rep:exercises}
% ------------------------------------------------------------

\begin{exercise}[$\star$]
Compute the TF-IDF vectors for the following corpus:
$d_1$ = ``the cat sleeps'', $d_2$ = ``the dog sleeps'',
$d_3$ = ``the cat and the dog''.
\end{exercise}

\begin{exercise}[$\star$]
Show that cosine similarity is invariant under positive scaling:
$\cos(\lambda \mathbf{u}, \mu \mathbf{v}) = \cos(\mathbf{u}, \mathbf{v})$ for
$\lambda, \mu > 0$.
\end{exercise}

\begin{exercise}[$\star\star$]
Implement a bigram language model with Laplace smoothing. Compute the
perplexity on a test corpus.
\end{exercise}

\begin{exercise}[$\star\star$]
Show that the Eckart-Young theorem implies that the truncated SVD gives the
best rank-$k$ approximation in the Frobenius norm.
\end{exercise}

\begin{exercise}[$\star\star\star$]
Compare the performance of TF-IDF + logistic regression, TF-IDF + SVM, and
LSA + logistic regression on sentiment classification (IMDB dataset). Analyze
the impact of the LSA dimension $k$.
\end{exercise}
